\section{Types and Metadata == Business Modeling}\label{sec:02}

  In recent years there was a spike in popularity of the Domain-Specific Languages (DSL) concept, which can be explainded by the programming issues outlined in the introduction, which affect not only the business oriented information systems, but also other domains dependent on various informations systems.

  The concept of DSL is not new.
  In fact, there are well established solutions for business software what exist for a long time and explictly utilise this concept.
  For example, SAP Business Suite offers a specialised language ABAP and Microsoft Dynamics Axapta provides X++.
  These are examples of so called \emph{external DSL}, which are separate or external to the programming language used natively for creation of the mentioned products.
  The main disadvantage with external DSLs is this: if there is a need to express a complex algorithm the expressive power of a DSL should be adequate to implement such algorithm, which in turn makes that DSL yet another general purpose programming language.
  
  An opposite to external DSL is the concept of \emph{interna DSL} -- a specialised language, which utilises the capabilities of a specific general purpose programming language of expressing language-like constructs that can be used together with the general purpose programming language.
  Internal DSLs share the development and execution infrastructure of the corresponding general purpose programming language taking a significant advantage of reusing debugging, profiling tools, code editors, existing programming libraries etc.
  
  Different general purpose programming languages provide different degrees of support for developing sophisticated internal DSLs.
  For example, scripting languages generally provide better ways of implementing internal DSLs (e.g. Ruby).
  However, in our strong oppinion statically type langugages are of essential importantce for constructing complex business oriented information systems.
  
  From the very beginning the Trident Genesis Platform was envisaged as an application platform founded on the concept of internal DSLs and domain-driven development.
  It was decided to tap into the power of modern IDEs and existing programming skills of software developers, which resulted in the choice of Java as the host general purpose programming language\footnote{Java does not provide the best support for developing internal DSLs (e.g. Scala is a lot more adequate), but the maturity of the existing Java infrastructure and development experience played a critical part in this decision.}.

%   \begin{figure}[!htp]
%     \centering
%     \includegraphics[width=10cm]{sections/01-intro/images/00-splash.jpg}
%     \caption{Example of an Embedded Image}\label{fig:01}
%   \end{figure}


%  // calculates the average yearly maintenance cost (for the last 3 years) per sectors of NORTH division. Only vehicles of MERCEDES make and models starting from 315 or 316 or VITO model are taken into account.
%         select(WorkOrder.class).
%         where(). //
%         prop("vehicle.model.make.key").eq().val("MERCEDES").and().
%         begin().prop("vehicle.model.key").starts_with().any_of_values("315", "316" ).or().prop("vehicle.model.key").eq().val("VITO").end().and().
%         year_of().prop("actualStart").in().values(2009, 2010, 2011).and().
%         prop("vehicle.station.zone.sector.division.key").eq().val("NORTH").
%         yield_and_group().prop("vehicle.station.zone.sector").as("sector").
%         yield().begin_expr().sum_of().prop("actualCost").div().val(3).end_expr().as("averageYearlyMaintenanceCostPerSector").model_as_aggregate();
 